\textbf{De fuerzas a aceleración y a distancia.} Se calcula primero $C_D=C_{D0}+kC_L^2=0.03428$. Para cada velocidad terrestre, $V_a=V_g+U$, $q=\rho V_a^2/2$, $L=qSC_L$ y $N=mg_0-L$; la fuerza de rueda utiliza esa reacción normal, no todo el peso. En frenado se define $b=-\dot V_g>0$ y $b=[D+\mu_bN]/m$.\[\fitmath{b(V_g)=2.8439285+(-8.244418605\times10^{-5})(V_g+0)^2\quad[\mathrm{m/s^2}]}\]\begin{center}\begin{tabular}{rrrrr}
\toprule $V_g$ & $V_a$ & $D$ (N) & $\mu N$ (N) & $b(V_g)$\\\midrule
$0$ & $0$ & $0$ & $4891.557$ & $2.843928$\\
$7.25$ & $7.25$ & $21.08156$ & $4863.022$ & $2.839595$\\
$14.5$ & $14.5$ & $84.32623$ & $4777.416$ & $2.826595$\\
$21.75$ & $21.75$ & $189.734$ & $4634.741$ & $2.804927$\\
$29$ & $29$ & $337.3049$ & $4434.995$ & $2.774593$\\
\bottomrule\end{tabular}\end{center}Como $ds/dt=V_g$ y $dV_g/dt$ es la aceleración con signo, se cambia la variable de integración. Con aceleración o desaceleración positiva en los límites elegidos:\[\fitmath{s=\int_0^{V_{g,f}}\frac{V_g}{b(V_g)}dV_g,\qquad t=\int_0^{V_{g,f}}\frac{1}{b(V_g)}dV_g}\]Para reproducir el cálculo a mano o con una hoja, Simpson compuesto usa un número par $n$ de intervalos, $h=V_{g,f}/n$ y $I\simeq(h/3)[f_0+f_n+4\sum_{i\ impar}f_i+2\sum_{i\ par,\ interior}f_i]$. Se aplica una vez con $f=V_g/b(V_g)$ y otra con $f=1/b(V_g)$. El refinamiento da:\begin{center}\begin{tabular}{rrr}
\toprule $n$ & $s$ (m) & $t$ (s)\\\midrule
$4$ & $149.6914$ & $10.28127$\\
$8$ & $149.6911$ & $10.28127$\\
$16$ & $149.6911$ & $10.28127$\\
$32$ & $149.6911$ & $10.28127$\\
$64$ & $149.6911$ & $10.28127$\\
\bottomrule\end{tabular}\end{center}\[\fitmath{s=\frac{1}{2a_2}\ln\!\left(\frac{a_0+a_2V_f^2}{a_0}\right)\approx149.6911\,\mathrm m}\]La primitiva sale de sustituir $u=a_0+a_2V^2$, $du=2a_2VdV$. Si $a_2\to0$, el límite es $V_f^2/(2a_0)$.La fase de demora se añade después: $s_d=28(1.2)=33.6$ m; el total terrestre es $s_b+s_d=183.2911\,\mathrm m$.
